% Front matter through the threshold-batch depth theorem.
\maketitle

\begin{abstract}
This note audits an obstacle/linear-complementarity route to the project's
second open problem (OP2).  In the shared normalization, regularized
PageRank is exactly the nonnegative quadratic
\[
 \min_{\vx\geq0}\ \phi(\vx)
 :=\tfrac12\vx^\trans\mQ\vx-\vc^\trans\vx,
 \qquad
 \vc=\vb-\alpha\rho\mD^{1/2}\one .
\]
Its KKT system is the symmetric nonsingular M-matrix LCP
$\vx\geq0$, $\vw=\mQ\vx-\vc\geq0$, and
$x_iw_i=0$.  This reduction has no sign or scaling slack.

The first positive ingredient is an exact safe-pivot theorem.
Starting from zero, solve the current principal face exactly and activate any
outside coordinate with negative complementarity slack.  Every activated
coordinate belongs to the unknown optimum support; after a batch activation,
the new exact face solution is positive and increases coordinatewise.  Thus
there are no false activations or deletions and at most $|\mathcal S_\rho^*|$ pivots.
Only graph-boundary coordinates can violate.  A local minimum-norm KKT
subgradient gives an objective-gap certificate in one charged terminal scan.

We further remove the hidden complementarity-margin issue.  After an
approximate face solve with residual norm $\delta$, it is enough to report a
boundary slack below the known threshold $-\delta/\alpha$.  Such a report is
provably negative for the exact face and hence support-safe.  If no such key
exists and
$\delta/\alpha\leq\sqrt{2\alpha\epsobj}/(1+2\sqrt{|\partial U|})$,
orthant projection already meets the objective target.  Thus the missing data
structure need not resolve all arbitrarily small signs.
Exact face changes also satisfy a global energy telescope: the sum of squared
$\mQ$-norm corrections, and hence the sum of squared full-slack changes, is at
most $\alpha$.  This magnitude bound does not by itself pay for communicating
dense corrections to boundary state.

The main result closes the repeated-face gap and proves OP2.  Order the unknown
optimal support by threshold-certified admission batches and take the block
Cholesky factor of its principal matrix.  Keeping only the diagonal and first
block subdiagonal produces a block-bidiagonal nonsingular M-matrix.  Its inverse
is entrywise no larger than the full inverse, so its smallest singular value is
at least $\sqrt\alpha$; its two-block row structure gives largest singular value
at most $\sqrt2$.  Chebyshev inverse decay on this block chain shows that the
unreleased face energy after $J$ batches is at most
$8q^{2J}+\tau^2/(\alpha\rho)$, where
$q=(\sqrt{2/\alpha}-1)/(\sqrt{2/\alpha}+1)$ and $\tau$ is the declared residual
threshold.  The second term pays every delayed or numerically ambiguous release
at once; no complementarity margin is assumed.

Choose $\tau=\Theta(\sqrt{\alpha\rho\epsobj})$ and
$J=O(\alpha^{-1/2}\log(1/\epsobj))$.  At each batch, a fresh nearly-linear SDD
solve of the exposed degree-coordinate face is accepted only after a fully
charged exact active-residual scan; capped independent retries provide the
declared success probability.  One boundary scan then either admits all
certified violations above threshold or stops.  Every admitted row lies in
$\mathcal S_\rho^*$, hence every phase costs $\softO(1/\rho)$ and all
preprocessing is local to the exposed face.  Orthant
projection of the last approximate face gives objective gap at most
$\epsobj$.  This yields a graph-uniform randomized high-probability algorithm
with fully charged work
$\softO(1/(\rho\sqrt\alpha)\log(1/\epsobj))$ and no supplied support, global
preprocessing, boundary oracle, or uncharged materialization.

The earlier repeated-prefix obstruction remains a warning against counting only
terminal volume, but it is bypassed because only
$O(\alpha^{-1/2}\log(1/\epsobj))$ complete faces are solved.  On a
promised endpoint-seeded path, append-only scalar $LDL^\trans$ elimination
tests each successive boundary key in constant arithmetic and materializes
the solution once, giving exact-real $O(\vol(\mathcal S_\rho^*))$ work without
global preprocessing.

Finally, a four-vertex rational graph refutes the proposed invariant
$0\leq\vx_k\leq\vx_\rho^*$ for ordinary or orthant-projected face CG.  CG
started from an exact old face stays strictly feasible and decreases the
quadratic, yet overshoots one optimum coordinate at iteration three.
Projection is inactive, so it cannot repair the failure.  This is a precise
algorithmic obstruction, not a lower bound for all accelerated active-set
methods.
\end{abstract}

\tableofcontents

\section{Scope and claim status}
\label{sec:scope}

% Canonical source-aligned PageRank/RPPR reference formulation.
%
% This fragment is intentionally heading-free at the section level so that the
% active paper and every standalone note can input it. It is a manuscript
% reference convention, not an implementation-wide residual decision.
% Primary source: Fountoulakis and Martinez-Rubio, arXiv:2602.21138v2,
% PDF p. 3, Section 3 and equation (RPPR).

\subsection{Common source-aligned PageRank and RPPR model}
\label{subsec:shared-source-aligned-problem}

All documents in this manuscript workspace use the following reference
language. It follows the plain italic notation of the comparison paper:
\(x,s,A,D,Q,I\) denote column vectors or matrices as determined by their
displayed domains. This is the scoped exception to the project's usual
bold-vector and bold-matrix typography. A note may impose a stronger
assumption after this block, but it must not redefine these objects.

Let \(G=(V,E)\) be a finite, simple, undirected, unweighted graph without
isolated vertices, let \(V=[n]:=\{1,\ldots,n\}\), and let
\(A\in\{0,1\}^{n\times n}\) be its symmetric adjacency matrix. For
\(i\in V\), write
\[
    \mathcal N(i):=\{j\in V:\{i,j\}\in E\},
    \qquad
    d_i:=|\mathcal N(i)|,
    \qquad
    D:=\diag(d_1,\ldots,d_n).
\]
For \(S\subseteq V\), define
\[
    \mathcal N(S):=\bigcup_{i\in S}\mathcal N(i),
    \qquad
    \partial S:=\mathcal N(S)\setminus S,
    \qquad
    \operatorname{Ext}(S):=V\setminus(S\cup\partial S),
\]
\begin{equation}
    \vol(S):=\sum_{i\in S}d_i.
    \label{eq:shared-volume}
\end{equation}
The support of a vector is
\(\supp(x):=\{i\in[n]:x_i\neq0\}\). Matrix inequalities use the Loewner
order, and all unqualified vector inequalities are coordinatewise.

The seed is a column distribution
\begin{equation}
    s\in\R_+^n,
    \qquad
    \inner{\one}{s}=1.
    \label{eq:shared-seed}
\end{equation}
The single-seed specialization is always written \(s=e_v\), where \(v\in V\)
is the seed vertex. Thus \(s\) is never used as a vertex label. For
\(\alpha\in(0,1]\), define
\begin{equation}
\begin{aligned}
    \mathcal L&:=I-D^{-1/2}AD^{-1/2},\\
    Q&:=\alpha I+\frac{1-\alpha}{2}\mathcal L
      =\frac{1+\alpha}{2}I
       -\frac{1-\alpha}{2}D^{-1/2}AD^{-1/2},\\
    b&:=\alpha D^{-1/2}s.
\end{aligned}
    \label{eq:shared-pagerank-matrices}
\end{equation}
Since \(0\preceq\mathcal L\preceq2I\),
\(\alpha I\preceq Q\preceq I\). The shared smooth PageRank quadratic is
\begin{equation}
    f(x):=\frac12\inner{x}{Qx}-\inner{b}{x},
    \qquad
    \nabla f(x)=Qx-b.
    \label{eq:shared-pagerank-objective}
\end{equation}
Its unique unregularized minimizer and associated lazy personalized PageRank
vector are
\begin{equation}
    x^0:=Q^{-1}b,
    \qquad
    \pi:=D^{1/2}x^0.
    \label{eq:shared-ppr-solution}
\end{equation}

For a regularization parameter \(\rho>0\), reserve \(g_\rho\) for the
nonsmooth weighted \(\ell_1\) term and define
\begin{equation}
    g_\rho(x):=\alpha\rho\norm{D^{1/2}x}_1,
    \qquad
    F_\rho(x):=f(x)+g_\rho(x),
    \qquad
    x^\star(\rho):=\argmin_{x\in\R^n}F_\rho(x).
    \label{eq:shared-rppr-objective}
\end{equation}
When \(\rho\) is fixed, \(x^\star\) abbreviates \(x^\star(\rho)\), never the
unregularized point \(x^0\). Write
\[
    S^\star(\rho):=\supp(x^\star(\rho)),
    \qquad
    I^\star(\rho):=[n]\setminus S^\star(\rho).
\]
The source records \(x^\star(\rho)\geq0\),
\(\vol(S^\star(\rho))\leq1/\rho\), and the coordinatewise conditions
\begin{equation}
    \nabla_i f(x^\star(\rho))
    \in
    \begin{cases}
        \{-\alpha\rho\sqrt{d_i}\}, & x_i^\star(\rho)>0,\\
        [-\alpha\rho\sqrt{d_i},0], & x_i^\star(\rho)=0.
    \end{cases}
    \label{eq:shared-rppr-kkt}
\end{equation}

The common computational unit is one repeated adjacency-list scan: scanning
coordinate \(i\) costs \(d_i\), and scanning a set \(S\) costs \(\vol(S)\).
Iteration count and wall-clock time do not replace this degree-weighted work
measure.

Finally, accuracy symbols remain distinct. The APPR threshold is
\(\epsappr\); \(\epsobj\) denotes an objective-gap target;
\(\epspg\) denotes the source experiment's proximal fixed-point tolerance; and
\(\epsppr\) may denote a document-scoped degree-normalized PPR error target.
No equality or conversion among these quantities is assumed. In particular,
this shared model does not resolve the repository-wide residual decision.


This note specializes to the canonical point seed $\vs=\eunit{v}$ and the
nonzero regime
\begin{equation}
 0<\rho<1/d_v.
 \label{eq:nonzero-regime}
\end{equation}
OP2 asks for a feasible returned vector $\widehat\vx$ satisfying
\begin{equation}
 F_\rho(\widehat\vx)-F_\rho(\vx_\rho^*)\leq\epsobj
 \label{eq:op2-gap}
\end{equation}
in fully charged graph work
$\softO(1/(\rho\sqrt\alpha))$, with only polylogarithmic dependence on
$1/\epsobj$.  The algorithm receives neither $\mathcal S_\rho^*$ nor free global
preprocessing.  Adjacency entries, degree queries, arithmetic, state
reads/writes, boundary and certificate scans, materialization, and output all
count.
The source variational formulation, nonnegativity, and support-volume theorem
are \citet[Section~4 and Theorems~1--2]{fountoulakis2019variational}; the
shared normalization is \citet[Section~3]{fountoulakis2026complexity}; and the
target complexity is the open question of
\citet[Section~3]{fountoulakis2022open}.

\begin{center}
\small
\begin{tabularx}{\textwidth}{@{}p{0.18\textwidth}X@{}}
\toprule
Status & Statement \\
\midrule
\statussource & The canonical RPPR optimum is nonnegative and has
$\vol(\mathcal S_\rho^*)\leq1/\rho$.  Wei--Yang give a true-support growing active-set
method but pay a repeated principal-system factor. \\
\statusproved & RPPR is exactly a symmetric nonsingular M-matrix LCP with
$\vc=\vb-\alpha\rho\mD^{1/2}\one$ and
$\vw=\mQ\vx-\vc$. \\
\statusproved & Exact negative-violation pivots from zero activate only
$\mathcal S_\rho^*$, require no deletions, and increase the exact face solution
coordinatewise. \\
\statusproved & On a supplied exact support, ordinary CG plus one final
orthant projection has work
$\bigO(\vol(\mathcal S_\rho^*)\alpha^{-1/2}\log(1/\epsobj))$. \\
\statusproved & A locally computable minimum-norm KKT subgradient certifies
objective gap at most $\norm{\vg}^2/(2\alpha)$. \\
\statusproved & A residual-calibrated threshold reporter either returns an
exactly support-safe pivot or certifies that one orthant projection meets the
objective target; no unknown sign margin is required. \\
\statusproved & Exact face corrections have total squared $\mQ$-energy and
full-slack motion at most $\alpha$ across the entire nested sequence. \\
\statusproved & Threshold-batch block Cholesky decay bounds the exact
unreleased face gap after $J$ nonempty batches by
$8q^{2J}+\tau^2/(\alpha\rho)$, with
$q=(\sqrt{2/\alpha}-1)/(\sqrt{2/\alpha}+1)$. \\
\statusproved & Fresh local SDD solves and complete active-row scans for
$O(\alpha^{-1/2}\log(1/\epsobj))$ threshold batches give a fully charged
$\softO(1/(\rho\sqrt\alpha))$ OP2 algorithm with polylogarithmic objective
accuracy dependence. \\
\statusproved & Promised endpoint-seeded paths have an append-only scalar
$LDL^\trans$ solver with exact-real $O(\vol(\mathcal S_\rho^*))$ total work. \\
\statusconditional & The stronger persistent active-edge contract in
\cref{def:active-edge-contract} would remove fresh per-batch rebuilding and
remains useful beyond the proved OP2 bound. \\
\statusrefuted & Terminal support volume alone pays neither repeated face
solves nor repeated boundary scans; endpoint paths give quadratic cumulative
prefix volume. \\
\statusrefuted & Ordinary face CG, and projected CG when its projection is
inactive, does not maintain $\vzero\leq\vx_k\leq\vx_\rho^*$ even for a
four-vertex graph M-matrix. \\
\statusopen & Deterministic finite-precision/bit-complexity and a persistent
changing-face implementation sharper than the proved fresh-batch solver. \\
\bottomrule
\end{tabularx}
\end{center}

The obstacle reduction, safe-pivot theorem, threshold-batch depth theorem,
and fully charged OP2 algorithm are promoted in condensed form to the active
manuscript.  The algorithmic theorem uses the shared exact-real algebraic word model and a randomized
nearly-linear SDD solver with declared failure probability.  It proves
polylogarithmic numerical-accuracy dependence, not coefficient-bit complexity
or a deterministic finite-precision implementation.

\section{Exact obstacle and LCP reduction}
\label{sec:reduction}

Define
\begin{equation}
 \vc:=\vb-\alpha\rho\mD^{1/2}\one,
 \qquad
 \phi(\vx):=\frac12\vx^\trans\mQ\vx-\vc^\trans\vx .
 \label{eq:c-and-phi}
\end{equation}
The off-diagonal entries of $\mQ$ are nonpositive, and
$\alpha\mI\preceq\mQ\preceq\mI$.

\begin{theorem}[Exact obstacle reduction]
\label{thm:reduction}
The RPPR problem is equivalent, with the same optimum value and unique
optimum, to
\begin{equation}
 \min_{\vx\geq\vzero}\phi(\vx).
 \label{eq:obstacle}
\end{equation}
Its KKT system is
\begin{equation}
 \vx\geq\vzero,
 \qquad
 \vw:=\mQ\vx-\vc\geq\vzero,
 \qquad
 x_iw_i=0\quad(i\in[n]).
 \label{eq:lcp}
\end{equation}
Moreover, $\mQ$ is a symmetric nonsingular M-matrix (a Stieltjes matrix) and
a K-matrix, so \cref{eq:lcp} has a unique solution.  In the regime
\cref{eq:nonzero-regime}, $v\in\mathcal S_\rho^*$.
\end{theorem}

\begin{proof}
For any $\vx$, replacing it by $|\vx|$ leaves the diagonal quadratic and
weighted $\ell_1$ terms unchanged.  Because $Q_{ij}\leq0$ for $i\neq j$,
the off-diagonal quadratic contribution cannot increase, and because
$\vb\geq\vzero$, $-\vb^\trans|\vx|\leq-\vb^\trans\vx$.  Thus
$F_\rho(|\vx|)\leq F_\rho(\vx)$, and uniqueness from strong convexity makes
$\vx_\rho^*\geq\vzero$.  On the nonnegative orthant,
\[
 \alpha\rho\norm{\mD^{1/2}\vx}_1
 =\alpha\rho\one^\trans\mD^{1/2}\vx,
\]
so $F_\rho(\vx)=\phi(\vx)$ there.  This proves \cref{eq:obstacle}.

At a lower bound $x_i\geq0$, stationarity and the normal cone give
$w_i=(\nabla\phi(\vx))_i\geq0$ when $x_i=0$ and $w_i=0$ when $x_i>0$.
These are exactly \cref{eq:lcp}; in particular the slack sign is positive,
not negative.  The matrix $\mQ$ is symmetric positive definite and has
nonpositive off-diagonals.  Hence it is Stieltjes.  Every principal minor is
positive, so it is also a P-matrix; together with the Z-sign pattern this is a
K-matrix.  Strict convexity, or the P-matrix LCP theorem, gives uniqueness.

Finally,
\[
 c_v=\frac{\alpha}{\sqrt{d_v}}(1-\rho d_v)>0.
\]
If $(\vx_\rho^*)_v=0$, then
$w_v=\sum_{j\neq v}Q_{vj}(\vx_\rho^*)_j-c_v<0$, contradicting
\cref{eq:lcp}.  Thus $v\in\mathcal S_\rho^*$.
\end{proof}

\begin{remark}[Degree-coordinate boundary key]
\label{rem:boundary-key}
Set $y_i=x_i/\sqrt{d_i}$.  Once the seed is active, every inactive
$j\neq v$ satisfies
\begin{equation}
 \frac{w_j}{\sqrt{d_j}}
 =\alpha\rho-\frac{1-\alpha}{2d_j}
   \sum_{i\in\mathcal N(j)\cap\supp(\vx)}y_i.
 \label{eq:boundary-key}
\end{equation}
If $j$ is not on the graph boundary of the support, the sum is empty and the
slack is strictly positive.  Detecting violations is therefore a dynamic
active-edge problem: admitted coordinates increase and monotonically lower
the keys of their inactive neighbors, but a dense old-face correction may
change every cut-edge contribution.
\end{remark}

\section{What a supplied active set gives}
\label{sec:supplied}

For $S\subseteq[n]$, let $\mQ_{SS}$ denote the principal matrix and extend
vectors on $S$ by zero outside $S$.  Principal interlacing gives
\begin{equation}
 \alpha\mI\preceq\mQ_{SS}\preceq\mI,
 \qquad
 \kappa(\mQ_{SS})\leq1/\alpha.
 \label{eq:principal-spectrum}
\end{equation}

\begin{theorem}[Supplied exact-support CG]
\label{thm:supplied-cg}
Suppose $S=\mathcal S_\rho^*$ is supplied.  Run ordinary exact-arithmetic CG on
\begin{equation}
 \mQ_{SS}\vy=\vc_S
 \label{eq:face-system}
\end{equation}
from any $\vy_0$, and after the last iteration return
$\widehat\vx_S=[\vy_k]_+$ and zero outside $S$.  If
$\Delta_S:=\tfrac12\norm{\vy_0-(\vx_\rho^*)_S}_{\mQ_{SS}}^2$, then
\begin{equation}
 F_\rho(\widehat\vx)-F_\rho(\vx_\rho^*)
 \leq \frac{4\Delta_S}{\alpha}
 \left(\frac{\sqrt\kappa-1}{\sqrt\kappa+1}\right)^{2k},
 \qquad \kappa:=\kappa(\mQ_{SS}).
 \label{eq:supplied-cg-rate}
\end{equation}
Consequently a supported-row implementation has charged work
\begin{equation}
 \bigO\!\left(
   \frac{\vol(\mathcal S_\rho^*)}{\sqrt\alpha}
   \log\frac{\Delta_S}{\alpha\epsobj}
 \right),
 \label{eq:supplied-work}
\end{equation}
including one final projection and output scan.
\end{theorem}

\begin{proof}
Every optimum coordinate on its exact support is positive, so
$\mQ_{SS}(\vx_\rho^*)_S=\vc_S$.  The standard CG energy estimate
\citep{hestenes1952methods} is
\[
 \norm{\vy_k-(\vx_\rho^*)_S}_{\mQ_{SS}}
 \leq2q^k\norm{\vy_0-(\vx_\rho^*)_S}_{\mQ_{SS}},
 \quad q=\frac{\sqrt\kappa-1}{\sqrt\kappa+1}.
\]
Euclidean orthant projection is nonexpansive relative to the positive vector
$(\vx_\rho^*)_S$.  Since $\mQ_{SS}\preceq\mI$ and
$\mQ_{SS}\succeq\alpha\mI$,
\begin{align*}
 \phi([\vy_k]_+)-\phi((\vx_\rho^*)_S)
 &=\tfrac12\norm{[\vy_k]_+-(\vx_\rho^*)_S}_{\mQ_{SS}}^2\\
 &\leq\tfrac12\norm{[\vy_k]_+-(\vx_\rho^*)_S}_2^2\\
 &\leq\frac1{2\alpha}
       \norm{\vy_k-(\vx_\rho^*)_S}_{\mQ_{SS}}^2.
\end{align*}
This proves \cref{eq:supplied-cg-rate}.  A principal matrix-vector product
scans only rows in $S$ and costs $O(\vol(S))$.  The estimate
$q\leq\exp(-2/(\sqrt\kappa+1))$ and \cref{eq:principal-spectrum} give the
iteration and work bound.
\end{proof}

For the zero start, no graph-dependent polynomial is hidden in the logarithm:
since $F_\rho(\vzero)=0$ and
$\min f\geq-\norm{\vb}^2/(2\alpha)\geq-\alpha/2$, the initial RPPR gap is at
most $\alpha/2$.  The ordinary $\sqrt\kappa=O(1/\sqrt\alpha)$ factor therefore
enters exactly through the supplied principal solve.  What is hidden by
\cref{eq:supplied-work} is the cost of finding and maintaining $S$.

\section{Safe exact active-edge discovery}
\label{sec:safe-discovery}

For $U\subseteq[n]$ with $\mQ_{UU}$ nonempty, define the exact face-stationary
vector
\begin{equation}
 x_i^U=
 \begin{cases}
  (\mQ_{UU}^{-1}\vc_U)_i,&i\in U,\\
  0,&i\notin U,
 \end{cases}
 \qquad
 \vw^U:=\mQ\vx^U-\vc.
 \label{eq:face-vector}
\end{equation}
The empty face means $\vx^\emptyset=\vzero$ and $\vw^\emptyset=-\vc$.
We call a face \emph{reachable} if it arises from the empty face by repeatedly
adding only coordinates with strictly negative current slack.

\begin{theorem}[Safe batched K-LCP pivots]
\label{thm:safe-pivots}
Let $U$ be reachable.  Then
\begin{equation}
 U\subseteq\mathcal S_\rho^*,
 \qquad
 \vzero<\vx^U_U\leq(\vx_\rho^*)_U.
 \label{eq:reachable-invariant}
\end{equation}
Every coordinate $j\notin U$ with $w_j^U<0$ belongs to
$\mathcal S_\rho^*\setminus U$.  If $J$ is any nonempty collection of such coordinates,
then $U':=U\cup J$ is reachable and
\begin{equation}
 \vx^{U'}_{U'}>\vzero,
 \qquad
 \vx^{U'}_U\geq\vx^U_U.
 \label{eq:batch-monotonicity}
\end{equation}
If no outside slack is negative, $\vx^U=\vx_\rho^*$.  Consequently this exact
outer loop makes at most $|\mathcal S_\rho^*|$ pivots, has nested support, and never
activates a false coordinate or deletes a coordinate.
\end{theorem}

\begin{proof}
The empty-face slack at the seed is $w_v^\emptyset=-c_v<0$, so the first
pivot admits $v\in\mathcal S_\rho^*$.  Its scalar principal solve is positive.

Assume the invariant holds at a reachable $U$.  Restricted stationarity of
$\vx^U$ and the optimum equations on $\mathcal S_\rho^*$ give
\begin{equation}
 \mQ_{UU}\bigl((\vx_\rho^*)_U-\vx^U_U\bigr)
 =-\mQ_{U,\mathcal S_\rho^*\setminus U}(\vx_\rho^*)_{\mathcal S_\rho^*\setminus U}
 \geq\vzero.
 \label{eq:face-comparison}
\end{equation}
The inverse of a principal Stieltjes matrix is nonnegative
\citep[Chapter~6]{berman1994nonnegative}, so
$\vx^U_U\leq(\vx_\rho^*)_U$.

Now take $j\notin\mathcal S_\rho^*$.  Since $(\vx_\rho^*)_j=0$ and $w_j^*\geq0$,
\begin{align}
 w_j^U-w_j^*
 &=\mQ_{jU}\bigl(\vx^U_U-(\vx_\rho^*)_U\bigr)
   -\mQ_{j,\mathcal S_\rho^*\setminus U}(\vx_\rho^*)_{\mathcal S_\rho^*\setminus U}
 \geq0.
 \label{eq:no-false-violation}
\end{align}
Both terms are nonnegative because the coordinate differences in the first
are nonpositive and all displayed off-diagonal matrix entries are
nonpositive.  Hence $w_j^U\geq0$; a negative violator must lie in $\mathcal S_\rho^*$.

For a batch $J$ of negative violators, block elimination in the enlarged
principal equations gives
\begin{align}
 \mH_J\vx^{U'}_J&=-\vw^U_J,
 &\mH_J&:=\mQ_{JJ}-\mQ_{JU}\mQ_{UU}^{-1}\mQ_{UJ},
 \label{eq:schur-pivot}\\
 \vx^{U'}_U-\vx^U_U&=-\mQ_{UU}^{-1}\mQ_{UJ}\vx^{U'}_J.
 \label{eq:old-correction}
\end{align}
The Schur complement $\mH_J$ is again Stieltjes, so
$\mH_J^{-1}\geq0$.  Since $-\vw_J^U>\vzero$,
$\vx_J^{U'}>\vzero$; \cref{eq:old-correction} and the Z-sign pattern make the
old correction nonnegative.  This proves \cref{eq:batch-monotonicity} and
closes the induction.  If there is no negative outside slack, inside
stationarity and positivity give complementarity and outside nonnegativity;
uniqueness in \cref{thm:reduction} identifies the optimum.
\end{proof}

\begin{corollary}[Boundary-only discovery]
\label{cor:boundary-only}
At every nonempty reachable face, a negative violation can occur only in
$\partial U$.  Exposing every newly admitted vertex's adjacency list once costs
$\vol(U)$, and creates at most $\vol(U)$ boundary-candidate incidences.  No
adjacency list of a false boundary candidate need be scanned.
\end{corollary}

\begin{proof}
For $j\notin U\cup\partial U$ with $j\neq v$, sparsity gives
$w_j^U=-c_j=\alpha\rho\sqrt{d_j}>0$.  The seed is already in every nonempty
reachable face.  The incidence counts follow by charging each exposed edge to
its endpoint in $U$.
\end{proof}

\begin{proposition}[Exact-face energy and slack-motion telescope]
\label{prop:motion-telescope}
Let $U_0\subset U_1\subset\cdots\subset U_T$ be reachable exact faces and
write $\vDelta_t=\vx^{U_{t+1}}-\vx^{U_t}$.  Then
\begin{align}
 \sum_{t=0}^{T-1}\norm{\vDelta_t}_{\mQ}^2
 &=2\bigl(\phi(\vx^{U_0})-\phi(\vx^{U_T})\bigr)
 \leq\alpha,
 \label{eq:energy-telescope}\\
 \sum_{t=0}^{T-1}\norm{\vw^{U_{t+1}}-\vw^{U_t}}_2^2
 &\leq\sum_{t=0}^{T-1}\norm{\vDelta_t}_{\mQ}^2.
 \label{eq:slack-motion-telescope}
\end{align}
The bound controls squared numerical motion, not the number of state cells or
cut edges touched by an implementation.
\end{proposition}

\begin{proof}
Both consecutive face vectors are supported on $U_{t+1}$, and
$\vx^{U_{t+1}}$ is stationary for the quadratic restricted to that face.
The exact quadratic identity therefore gives
\[
 \phi(\vx^{U_t})-\phi(\vx^{U_{t+1}})
 =\tfrac12\norm{\vDelta_t}_{\mQ}^2.
\]
Summing telescopes.  Since every restricted minimum has objective at most
$\phi(\vzero)=0$ and at least the global minimum, the zero-start gap bound
following \cref{thm:supplied-cg} gives the final inequality in
\cref{eq:energy-telescope}.  Finally,
$\vw^{U_{t+1}}-\vw^{U_t}=\mQ\vDelta_t$ and
$\mQ^2\preceq\mQ$ because $\vzero\prec\mQ\preceq\mI$, proving
\cref{eq:slack-motion-telescope}.
\end{proof}

\begin{algorithm}[t]
\caption{Exact active-edge LCP skeleton}
\label{alg:skeleton}
\begin{algorithmic}[1]
\STATE $U\gets\emptyset$, $\vx^U\gets\vzero$.
\WHILE{a certified reporter returns $j\notin U$ with $w_j^U<0$}
  \STATE Choose a nonempty reported batch $J\subseteq\{j:w_j^U<0\}$.
  \STATE Continue the principal state from $U$ to $U\cup J$ and obtain
         $\vx^{U\cup J}$.
  \STATE Expose the adjacency lists of newly admitted vertices and update
         boundary keys.
  \STATE $U\gets U\cup J$.
\ENDWHILE
\STATE Materialize a feasible approximation and certify \cref{eq:op2-gap}.
\end{algorithmic}
\end{algorithm}

The theorem validates the outer logic of \cref{alg:skeleton}.  By itself it
says nothing about how cheaply an exact or sufficiently accurate principal
state can be continued, how a negative boundary key is located after a dense
old correction, or how ambiguous numerical signs are resolved.  Sections
\ref{sec:batch-depth}--\ref{sec:op2-algorithm} close those requirements at a
fixed threshold and objective accuracy.

\section{A local objective-gap certificate}
\label{sec:certificate}

For a feasible $\vx\geq\vzero$, set $\vw=\mQ\vx-\vc$ and define
\begin{equation}
 g_i(\vx):=
 \begin{cases}
  w_i,&x_i>0,\\
  \min\{w_i,0\},&x_i=0.
 \end{cases}
 \label{eq:minimal-subgradient}
\end{equation}

\begin{theorem}[KKT subgradient certificate]
\label{thm:kkt-certificate}
The vector $\vg(\vx)$ is the minimum-Euclidean-norm element of the
subdifferential of $\phi+I_{\R_+^n}$ at $\vx$, and
\begin{equation}
 F_\rho(\vx)-F_\rho(\vx_\rho^*)
 \leq\frac{\norm{\vg(\vx)}_2^2}{2\alpha}.
 \label{eq:gap-certificate}
\end{equation}
Thus $\norm{\vg(\vx)}_2\leq\sqrt{2\alpha\epsobj}$ is a sufficient OP2
stopping rule.  If $\vx$ is stored on $U$ containing the seed, one scan of
the rows in $U$ computes all positive-support residuals and all possibly
negative outside entries, at work $O(\vol(U))$ plus $O(|U|)$ state and output
words.
\end{theorem}

\begin{proof}
At $x_i>0$ the normal cone is zero.  At $x_i=0$ it is $(-\infty,0]$, so the
smallest magnitude in $w_i+(-\infty,0]$ is zero when $w_i\geq0$ and $w_i$
when $w_i<0$.  This proves \cref{eq:minimal-subgradient}.  Strong convexity
gives, for every subgradient $\vg$,
\[
 (\phi+I)(\vy)\geq(\phi+I)(\vx)+\vg^\trans(\vy-\vx)
       +\frac\alpha2\norm{\vy-\vx}_2^2.
\]
Minimizing the right side over $\vy\in\R^n$ and taking
$\vy=\vx_\rho^*$ yields \cref{eq:gap-certificate}.  For locality, a sparse
row scan accumulates $\mQ\vx$ on $U\cup\partial U$.  Every farther nonseed
coordinate has positive slack by \cref{rem:boundary-key}, so its minimum-norm
subgradient entry is zero.
\end{proof}

The certificate explicitly charges the terminal scan.  It also avoids a
minimum positive coordinate or strict-complementarity assumption.  Combined
with the shared source conversion, objective gap $\epsobj$ yields the usual
degree-normalized PageRank accuracy
$\rho+\sqrt{2\epsobj/\alpha}$; this is a consequence, not a replacement for
the OP2 objective target.

The next result removes the different margin that arises when an approximate
face solution is used to decide whether a boundary slack is negative.

\begin{theorem}[Margin-free approximate-face dichotomy]
\label{thm:threshold-dichotomy}
Let $U$ be a nonempty reachable face, let $B=\partial U$, and put
$b_U=|B|$.  Let $\vz$ be any vector supported on $U$ and define
\begin{equation}
 \vr_U:=\mQ_{UU}\vz_U-\vc_U,
 \qquad
 \delta:=\norm{\vr_U}_2,
 \qquad
 \eta:=\delta/\alpha,
 \qquad
 \tau:=\sqrt{2\alpha\epsobj}.
 \label{eq:threshold-parameters}
\end{equation}
Suppose
\begin{equation}
 \eta\leq \frac{\tau}{1+2\sqrt{b_U}}.
 \label{eq:face-residual-target}
\end{equation}
For $j\in B$, form the approximate boundary slack
$\widehat w_j=(\mQ\vz-\vc)_j$.  Then the following dichotomy provides a
valid action.
\begin{enumerate}[label=(\roman*)]
\item If some $j\in B$ satisfies
\begin{equation}
 \widehat w_j<-\eta,
 \label{eq:safe-threshold-report}
\end{equation}
then the exact face slack satisfies $w_j^U<0$.  Thus $j$ is a support-safe
pivot covered by \cref{thm:safe-pivots}.
\item If no boundary coordinate satisfies \cref{eq:safe-threshold-report},
then the feasible vector
\begin{equation}
 \widehat\vx_U=[\vz_U]_+,
 \qquad
 \widehat\vx_{\bar U}=\vzero
 \label{eq:threshold-output}
\end{equation}
satisfies
$F_\rho(\widehat\vx)-F_\rho(\vx_\rho^*)\leq\epsobj$.
\end{enumerate}
Consequently the reporter need only return a key below the known,
objective-derived threshold $-\delta/\alpha$, or certify that no such key
exists.  It never has to resolve an unknown minimum nonzero complementarity
slack.
\end{theorem}

\begin{proof}
Let $\vx^U$ be the exact face vector and set
$\ve_0=\vz-\vx^U$.  Restricted stationarity gives
$\mQ_{UU}(\ve_0)_U=\vr_U$, so
\begin{equation}
 \norm{\ve_0}_2\leq\delta/\alpha=\eta.
 \label{eq:face-error-bound}
\end{equation}
Because $\norm{\mQ}_2\leq1$, every boundary slack error obeys
$|\widehat w_j-w_j^U|\leq\eta$.  Hence
\cref{eq:safe-threshold-report} implies $w_j^U<0$, proving the first case.

Suppose there is no report.  Then
$w_j^U\geq\widehat w_j-\eta\geq-2\eta$ for every $j\in B$.
Only boundary coordinates can have negative exact slack, so, with the negative
part extended by zero,
\begin{equation}
 \norm{[\vw^U_{\bar U}]_-}_2\leq2\sqrt{b_U}\,\eta.
 \label{eq:exact-face-negative-mass}
\end{equation}
Let $\ve=\widehat\vx-\vx^U$.  Orthant projection is nonexpansive relative to
the nonnegative exact face vector, and therefore
$\norm{\ve}_2\leq\norm{\ve_0}_2\leq\eta$.  Define
\begin{equation}
 \vh:=\mQ\ve+
 \begin{bmatrix}
  \vzero_U\\ [\vw^U_{\bar U}]_-
 \end{bmatrix}.
 \label{eq:constructed-subgradient}
\end{equation}
At positive coordinates of $\widehat\vx$, $\vh$ equals the gradient.  At a
zero coordinate in $U$, choosing zero normal makes the same identity valid;
outside $U$, subtracting the positive part of $\vw^U$ is an admissible
nonpositive normal.  Thus
$\vh\in\partial(\phi+I_{\R_+^n})(\widehat\vx)$.  Equations
\cref{eq:face-residual-target,eq:exact-face-negative-mass} give
\[
 \norm{\vh}_2
 \leq\norm{\mQ\ve}_2+\norm{[\vw^U_{\bar U}]_-}_2
 \leq(1+2\sqrt{b_U})\eta
 \leq\tau.
\]
The strong-convexity argument in \cref{thm:kkt-certificate}, which applies to
every subgradient and not only the minimum-norm one, proves the objective
claim.
\end{proof}

\begin{corollary}[Overlapping thresholds remove finite-precision sign margins]
\label{cor:interval-threshold}
In the setting of \cref{thm:threshold-dichotomy}, let
$\bar\eta\geq\delta/\alpha$ be a certified error bound with
\begin{equation}
 \bar\eta\leq\frac{\tau}{1+4\sqrt{b_U}}.
 \label{eq:interval-residual-target}
\end{equation}
For each $j\in B$, suppose a certified interval $I_j$ contains
$\widehat w_j$ and has width strictly less than $2\bar\eta$.  If
$\sup I_j<-\bar\eta$ for some $j$, then $w_j^U<0$ and the pivot is safe.  If
there is no such report, the projected vector in
\cref{eq:threshold-output} has objective gap at most $\epsobj$.
Thus every key need only be refined to the predetermined additive scale
$\bar\eta$; no key-dependent separation margin is required.
\end{corollary}

\begin{proof}
A report gives
$w_j^U\leq\widehat w_j+\bar\eta<0$.  Otherwise
$\sup I_j\geq-\bar\eta$ for every $j$.  Since the interval width is less than
$2\bar\eta$, its lower endpoint is greater than $-3\bar\eta$, and hence
$\widehat w_j>-3\bar\eta$.  Therefore
$w_j^U>-4\bar\eta$ and
$\norm{[\vw^U_{\bar U}]_-}_2\leq4\sqrt{b_U}\bar\eta$.  Repeating the
constructed-subgradient proof with
$\norm{\widehat\vx-\vx^U}_2\leq\delta/\alpha\leq\bar\eta$ gives a
subgradient of norm at most
$(1+4\sqrt{b_U})\bar\eta\leq\tau$.
\end{proof}

For a fixed supplied face, ordinary CG reaches
\cref{eq:face-residual-target} in
$O(\sqrt{\kappa(\mQ_{UU})}\log(b_U/(\alpha\epsobj)))$ iterations up to
constant and initial-gap factors.  The dependence on $b_U$ and
$1/\epsobj$ is logarithmic.  This theorem closes sign robustness, but it does
not amortize the matrix-vector products or threshold reports across changing
faces.

\section{Threshold-batch Cholesky decay}
\label{sec:batch-depth}

The preceding paragraph treats every face independently.  The next argument
instead orders the unknown optimal support by the batches in which it is
admitted.  Its central fact is that delayed releases form a well-conditioned
block-bidiagonal Cholesky chain.  This gives the square-root depth bound even
though ordinary Bellman contraction gives only the factor $1-\Theta(\alpha)$.

Write the outside face residual as
\begin{equation}
 \vr^U:=\vc-\mQ\vx^U=-\vw^U.
 \label{eq:outside-face-residual}
\end{equation}
Fix $\vartheta\geq0$.  A nested sequence
$U_k=B_0\mathbin{\dot\cup}\cdots\mathbin{\dot\cup}B_k$ is called
\emph{$\vartheta$-batched} if $B_0=\{v\}$, every later $B_k$ is a nonempty
set of coordinates with $r_i^{U_{k-1}}>0$, and, immediately before admitting
$B_k$, every coordinate that will be admitted after $B_k$ has residual at
most $\vartheta$.  Coordinates outside the seed have empty-face residual
$c_i=-\alpha\rho\sqrt{d_i}<0$, which supplies the analogous bound before
$B_0$ without adding it to the definition.  Any finite such sequence may be
extended to a partition of $\mathcal S_\rho^*$: after its last declared batch,
repeatedly admit all exact positive residuals.  After each such all-positive
batch is selected, every still-future coordinate has current residual at most
zero, so the extension remains $\vartheta$-batched.  For the first exact
continuation block, the different two-face-back causal bound used below comes
from the last declared batch condition; its newly positive residual on the
immediately preceding face need not be small.

\begin{theorem}[Threshold-batch energy-depth bound]
\label{thm:batch-depth}
Let $U_0\subset U_1\subset\cdots$ be a $\vartheta$-batched reachable
sequence, extended as above to the full optimal support.  Put
\begin{equation}
 q_\alpha:=
 \frac{\sqrt{2/\alpha}-1}{\sqrt{2/\alpha}+1}.
 \label{eq:batch-q}
\end{equation}
Then for every included face $U_J$,
\begin{equation}
 \phi(\vx^{U_J})-\phi(\vx_\rho^*)
 \leq 8q_\alpha^{2J}+\frac{\vartheta^2}{\alpha\rho}.
 \label{eq:batch-depth-gap}
\end{equation}
In particular, a constant reduction needs only
$O(\alpha^{-1/2})$ batches up to logarithmic factors; no terminal-support
cardinality, graph diameter, or complementarity margin occurs in the depth,
provided the requested gap lies above the displayed additive threshold floor.
\end{theorem}

\begin{proof}
Restrict $\mQ$ and $\vc$ to $S:=\mathcal S_\rho^*$ and order $S$ by the
blocks $B_0,B_1,\ldots,B_K$.  Block elimination gives
\begin{equation}
 \mQ_{SS}=\mL\mDelta\mL^\trans,
 \qquad
 \mDelta=\operatorname{diag}(\mDelta_0,\ldots,\mDelta_K),
 \qquad
 \vg=\mL^{-1}\vc_S.
 \label{eq:block-ldl}
\end{equation}
Here $\mL$ is unit block lower triangular, every pivot
$\mDelta_k$ is Stieltjes, and $\vg_k$ is exactly the positive Schur residual
on the batch $B_k$ when it is admitted.  If
$\mDelta_k=\mR_k\mR_k^\trans$ is its lower Cholesky factor and
\begin{equation}
 \vz_k:=\mR_k^{-1}\vg_k,
 \label{eq:block-z}
\end{equation}
then $\vz_k\geq\vzero$ because $\mR_k$ is a triangular nonsingular
M-matrix.  Standard completion of squares in block elimination gives the
exact face-gap identity
\begin{equation}
 2\bigl(\phi(\vx^{U_J})-\phi(\vx_\rho^*)\bigr)
 =\sum_{k>J}\norm{\vz_k}_2^2.
 \label{eq:block-gap-identity}
\end{equation}

Put $\mC:=\mL\operatorname{diag}(\mR_0,\ldots,\mR_K)$, so
$\mQ_{SS}=\mC\mC^\trans$.  This is the unique scalar lower Cholesky factor
under the block-respecting coordinate order.  Its recursion
$C_{ij}=(Q_{ij}-\sum_{\ell<j}C_{i\ell}C_{j\ell})/C_{jj}$ shows inductively
that every off-diagonal entry is nonpositive.  Let
$\widehat{\mC}$ retain from $\mC$ only its diagonal blocks and its first block
subdiagonal.  Both $\mC$ and $\widehat{\mC}$ are triangular nonsingular
M-matrices, and $\mC\leq\widehat{\mC}$ entrywise.  Inverse monotonicity and
nonnegativity therefore give
\begin{equation}
 \vzero\leq\widehat{\mC}^{-1}\leq\mC^{-1},
 \qquad
 \norm{\widehat{\mC}^{-1}}_2
 \leq\norm{\mC^{-1}}_2\leq\frac1{\sqrt\alpha}.
 \label{eq:bidiagonal-inverse-bound}
\end{equation}
Indeed, the inverse ordering follows directly from
$\mC^{-1}-\widehat{\mC}^{-1}
=\mC^{-1}(\widehat{\mC}-\mC)\widehat{\mC}^{-1}\geq0$.
The passage from entrywise to spectral norm is valid because both inverses
are nonnegative: $|\widehat{\mC}^{-1}z|\leq
\mC^{-1}|z|$ for every $z$.

There is also a graph-independent upper bound.  The $k$th block row of
$\mC$ has squared row-map matrix
\begin{equation}
 \sum_{i\leq k}\mC_{ki}\mC_{ki}^\trans
 =\mQ_{B_kB_k}\preceq\mI.
 \label{eq:block-row-bound}
\end{equation}
Deleting all but two input blocks cannot increase this row-map norm.  Each
input block occurs in at most two retained rows, whence
\begin{equation}
 \norm{\widehat{\mC}}_2\leq\sqrt2.
 \label{eq:bidiagonal-upper-bound}
\end{equation}
Consequently
$\mA:=\widehat{\mC}\widehat{\mC}^\trans$ is block tridiagonal and
\begin{equation}
 \alpha\mI\preceq\mA\preceq2\mI.
 \label{eq:block-chain-spectrum}
\end{equation}

We next encode the batch chronology.  For $k\geq2$, immediately before
$B_{k-1}$ is admitted, the still-future block $B_k$ has residual at most
$\vartheta\one$; this is its residual after eliminating through $B_{k-2}$.
For $k=1$, the same bound follows from
$c_i=-\alpha\rho\sqrt{d_i}<0$ off the seed.  Eliminating $B_{k-1}$ subtracts
$\mL_{k,k-1}\vg_{k-1}$.  Hence, for every $k\geq1$,
\begin{equation}
 \vg_k+\mL_{k,k-1}\vg_{k-1}\leq\vartheta\one.
 \label{eq:causal-pivot-bound}
\end{equation}
Using $\mC_{k,k-1}=\mL_{k,k-1}\mR_{k-1}$, equations
\cref{eq:block-z,eq:causal-pivot-bound} say
\begin{equation}
 (\widehat{\mC}\vz)_0=\vg_0,
 \qquad
 (\widehat{\mC}\vz)_k\leq\vartheta\one\quad(k\geq1).
 \label{eq:chain-forcing}
\end{equation}
Because $\widehat{\mC}^{-1}\geq0$,
\begin{equation}
 \vzero\leq\vz
 \leq\widehat{\mC}^{-1}
 \begin{bmatrix}\vg_0\\ \vartheta\one_{S\setminus B_0}\end{bmatrix}.
 \label{eq:chain-domination}
\end{equation}

It remains to bound the two sources.  For the distributed threshold source,
\cref{eq:bidiagonal-inverse-bound} and $|S|\leq\vol(S)\leq1/\rho$ give
\begin{equation}
 \norm{\widehat{\mC}^{-1}
 [\vzero;\vartheta\one]}_2^2
 \leq\frac{\vartheta^2|S|}{\alpha}
 \leq\frac{\vartheta^2}{\alpha\rho}.
 \label{eq:threshold-forcing-bound}
\end{equation}

For the seed-block source and $J\geq1$, use a degree-$(J-1)$ Chebyshev
polynomial $p_{J-1}$ for $1/t$ on $[\alpha,2]$.  Its standard uniform error is
\begin{equation}
 \norm{\mA^{-1}-p_{J-1}(\mA)}_2
 \leq\frac2\alpha q_\alpha^J.
 \label{eq:inverse-chebyshev}
\end{equation}
Since $\mA$ is block tridiagonal and $\vg_0$ is supported on block zero,
$p_{J-1}(\mA)\vg_0$ is supported on blocks through $J-1$.  Multiplication by
the block upper-bidiagonal $\widehat{\mC}^\trans$ does not increase the
largest block index.  Therefore, writing $P_{>J}$ for the block-tail
projection and using
$\widehat{\mC}^{-1}=\widehat{\mC}^\trans\mA^{-1}$,
\begin{align}
 \norm{P_{>J}\widehat{\mC}^{-1}[\vg_0;\vzero]}_2
 &\leq\norm{\widehat{\mC}}_2
       \norm{\mA^{-1}-p_{J-1}(\mA)}_2\norm{\vg_0}_2\notag\\
 &\leq\frac{2\sqrt2}{\alpha}q_\alpha^J\norm{\vg_0}_2
 \leq2\sqrt2q_\alpha^J.
 \label{eq:seed-tail-bound}
\end{align}
The last inequality uses
$\vg_0=c_v=\alpha(1-\rho d_v)/\sqrt{d_v}\leq\alpha$.
For $J=0$, the same displayed tail bound follows directly from
$\norm{\widehat{\mC}^{-1}}_2\leq1/\sqrt\alpha$,
$\norm{\vg_0}_2\leq\alpha$, and $\alpha\leq1$.
Combining \cref{eq:chain-domination,eq:threshold-forcing-bound,eq:seed-tail-bound}
with $\norm{a+b}_2^2/2\leq\norm a_2^2+\norm b_2^2$ and then
\cref{eq:block-gap-identity} proves \cref{eq:batch-depth-gap}.
\end{proof}

The proof explains both known extremes.  At $\vartheta=0$ it recovers an
accelerated exact-batch depth bound without the earlier $\rho=0$ Krylov
containment.  On a long endpoint path the retained block-bidiagonal factor is
the whole scalar factor, and the decay scale is sharp up to constants.  A
delayed same-boundary release contributes to the distributed forcing only
once, at the block where that vertex is eventually admitted; it cannot create
an uncharged sequence of margin errors.
